\section{Action Laws and Their Limits}

\begin{law}[Observational idempotence of add]
\label{law:add-observational}
For valid $d$,
\[
  \Add{d}(\Add{d}(S))\HAequiv\Add{d}(S).
\]
The allow list contains one copy of $d$, and the active hosts file contains no
line with an active $d$ token.
\end{law}

\begin{boundary}[Add is not fully idempotent]
\label{boundary:add-backup}
Under full state equality,
\[
  \Add{d}(\Add{d}(S))\ne\Add{d}(S)
\]
in general.  After the first call, $B=\Some{S.H}$.  After the second call,
$B=\Some{\Remove{d}(S.H)}$.

The behavior test for repeated addition proves idempotence of the allow list
and active hosts result.  It does not require preservation of the first backup.
The distinction is semantic idempotence versus backup-history idempotence.
\end{boundary}

\begin{law}[Setwise commutativity of distinct adds]
\label{law:add-commute}
For valid, distinct $d$ and $e$,
\[
  \Add{d}(\Add{e}(S))\Setequiv\Add{e}(\Add{d}(S)).
\]
Full equality can fail because appended allow-list order differs and because
\path{hosts.bak} records the hosts file immediately before the final add.
\end{law}

\begin{law}[Restore idempotence]
\label{law:restore-idempotent}
If the original backup remains present and unchanged, and both confirmations
succeed, then
\[
  \Restore(\Restore(S))=\Restore(S).
\]
The model omits timestamps that can differ between physical copies.
\end{law}

\begin{law}[Prep-restore round trip]
\label{law:roundtrip}
If prep succeeds and $O$ is not modified before restore, then
\[
  \Restore(\Prep_E(S)).H=S.H.
\]
This is a round-trip law for active contents, not necessarily for the complete
state.  Cache state and physical timestamps can differ.
\end{law}

\begin{boundary}[Prep is not idempotent]
\label{boundary:prep}
In general,
\[
  \Prep_E(\Prep_E(S))\ne\Prep_E(S).
\]
A confirmed second preparation overwrites \path{hosts-ORIG} with the current
hblock-generated hosts file.  That changes the restoration point.  The
confirmation prompt is therefore semantically important rather than merely a
user-interface convenience.
\end{boundary}

\begin{law}[Verbose noninterference]
\label{law:verbose}
Let $q(c)$ and $v(c)$ differ only in verbosity.  For any successful command,
\[
\StateOf(\Exec(q(c),S))=\StateOf(\Exec(v(c),S)).
\]
Only the trace is expected to differ:
\[
\TraceOf(\Exec(q(c),S))\ne\TraceOf(\Exec(v(c),S)).
\]
The model ignores incidental access-time changes a filesystem might make while
objects are inspected.
\end{law}

\begin{law}[Abstract idempotence of flush]
\label{law:flush}
\[
  \Flush(\Flush(S)).K=\Fresh.
\]
Process identifiers, timing, and diagnostics may differ, so the law applies
only to the abstract cache state.
\end{law}
